\section{Introduction}\label{sec:current-overview}\label{sec:introduction}

A predictive state records the information in an experimental past that
can affect future laws.  Its linear dimension measures one form of
compression.  A device storing a probability distribution on finitely
many labels faces additional constraints: its transition and readout
tables must be positive, its initialization must be implemented, and a
single device may have to serve several preparations and later requests.
These constraints are also present when the aim is to implement an
optimal experimental policy rather than to reproduce a prescribed law.
The resulting resource can be much larger than the predictive dimension
or the resource of any one preparation considered separately.

The question of this paper is how predictive, convex and spectral
structure controls such finite positive implementations.  We work
within a theory of prepared causal experiments, acquired task laws and
finite-resource statistical decisions.  A terminal response assigns a
probability row to each command word.  An instrument response describes
joint command--report histories and the state on which a subsequent
terminal test acts.  All controllers have one fixed parameter-independent
program.  Persistent labels, policy storage, physical state, program
bits and arithmetic work are different resources.  An initializer is
counted when it is a pure label of the device; freely prepared mixed
starts are always stated as a separate convention.

Three structural results organize the current article.  First, we
classify the common positive realization of multiple preparations under
finite-group response constraints.  The exact initialization term is
geometric, and the joint rather than individual Fourier multiplicity
governs small-contrast memory.  Second, we identify the exact number of
fixed command rows needed for robust optimal control of finite reset
experiments with tied actions.  This number is a finite transversal
invariant coupled to the optimal initial row.  Third, we give finite
operational memory certificates which do not require physical density
matrices or Born readouts at hidden simulator labels.  The original
algebraic qubit experiment then has a provably diverging approximate
memory requirement even in this larger decoder class.

\subsection*{One controller for several preparations}

Fix a finite group $G$, a finite command alphabet whose words map onto it, and
a finite family of preparation responses $b_i:G\to\Delta(C)$.
Every preparation uses the same command matrices and terminal decoder.
The decoder does not receive the preparation identity as an extra
input.  For closed convex specifications $K_{i,g}$, imposed on every
word with image $g$, Theorem~\ref{thm:r77-shared-normal-form} gives
a common permutation action $P_g$ on a finite set $X$, one decoder
$D$, and preparation rows $\lambda_i$ satisfying
\[
 \lambda_iP_gD\in K_{i,g}.
\]
Its pure-start label charge is exactly
$|X|+\iota(\Lambda)$, where $\iota(\Lambda)$ counts the distinct
nonpoint rows among the $\lambda_i$.  In the freely mixed convention
it is $|X|$.  The original stochastic controller need not satisfy the
group relations, and the resulting action need not be faithful.

For exact targets let $\mathcal B$ be the set of distinct prescribed
response tables.  Consider all $G$-invariant polytopes $P$ in the
product simplex of probability-valued functions on $G$ that contain
$\mathcal B$.  With $v(P)$ the number of vertices,
Theorem~\ref{thm:r77-shared-polytope} proves
\begin{equation}\label{eq:r77-intro-polytope}
 \kappa_G(\mathcal B)=\min_P v(P),\qquad
 m_G(\mathcal B)=\min_P\{v(P)+|\mathcal B\setminus\operatorname{vert}P|\}.
\end{equation}
The minimum polytope is allowed to extend beyond the target linear
span.  Its vertices are continuing permutation states.  Every distinct
prescribed nonvertex response has one counted transient initializer.
The proof also accounts for nonextreme continuing labels in an
unreduced controller, so the formula does not charge such labels twice.

Let $V_{\mathcal B}$ be the real span of all translated preparation
responses.  If $A_\rho(b_i)$ is the stacked Fourier block for an
irreducible representation $\rho$, its multiplicity in the joint
response module is
\[
 r_\rho(\mathcal B)
 =\rank[\,A_\rho(b_1)\ \cdots\ A_\rho(b_{|I|})\,].
\]
This rank can exceed the degree of $\rho$, a possibility excluded by
a single cyclic preparation.  Lemma~\ref{lem:r77-joint-multiplicity}
and Theorem~\ref{thm:r77-shared-low-contrast} show that when all
preparations have a common group average $\pi$, the exact small-positive-
contrast mixed memory of $\pi+t(b_i-\pi)$ is the minimum permutation
degree containing this joint module.  It is an integer cover by the
fixed-space multiplicities of subgroups.  A common positive decoder
is constructed explicitly.  The common-mean assumption is substantive:
a four-preparation trivial-group example has linear rank three and
positive memory four.

Theorem~\ref{thm:r77-joint-robust} replaces exact joint ranks by singular
values above the specified approximation error, giving robust resource
certificates.  Two exact families expose the separate costs.  With
$q$ preparations on the fixed group $C_2$, each preparation alone
needs two labels but a shared controller needs $2q$; the obstruction
persists at a positive explicitly bounded error.  A second family has
shared mixed memory two and pure memory $r+1$ for $r$ preparations.
These are growing finite experiments, with arbitrary finite preparation
count, rather than a change of numerical prior in one fixed example.

The inherited finite-group convex-specification theorem and its
attained all-word minimax principle remain in the article with their
full proofs.  The supplement retains the exact one-preparation normal
form, the complete mean--contrast law for the displayed family on
every $S_n$, $n\ge5$, its boundary rigidity, normal-stabilizer
convexity and quotient duals.  The shared-device results extend those
premises without replacing a general positive-realization question
by linear dimension alone.

\subsection*{Optimal feedback with ties and a counted initial row}

The finite reset instrument has failure branches
$\mathcal I_{a,i}(\rho)=\tr(E_{a,i}\rho)\sigma_i$ and a success
branch on which the experiment must terminate.  Its observed report
prepares a specified state $\sigma_i$; no external reset or state
oracle is supplied.  Commands are arbitrary in a finite set, call
costs have a positive floor, and the initial density matrix is unknown.
The inherited primitive theorem reduces its unrestricted optimum to a
finite Bellman vector $v$ and the spectral minimization
\[
 V_* =\min_{p\in\Delta(\mathcal A)}
       \lambda_{\max}\Bigl(\sum_a p_aL_a\Bigr),
 \qquad L_a=G_a+\sum_i v_iE_{a,i}.
\]
Its preparation saddle supplies a residual identity, but did not itself
classify the exact number of rows of every optimal tied-action controller.

Assume now that every failure effect is positive definite and there
is no closed failure set.  Every reset report then occurs with positive
probability on the first call under a single least-favorable preparation,
including a singular such preparation.  Let $\mathcal A_i^*$ be the
set of Bellman-optimal actions at reset state $i$, $\mathcal P_*$ the
set of optimal initial rows, and $I(p)=\{i:\supp p\subseteq\mathcal
A_i^*\}$.  If $t(J)$ denotes the least number of actions meeting
every $\mathcal A_i^*$ for $i\in J$, then
Theorem~\ref{thm:r77-reset-cover} proves
\begin{equation}\label{eq:r77-intro-reset}
 m_*=1+\min_{p\in\mathcal P_*}t(\{1,\ldots,s\}\setminus I(p)).
\end{equation}
There is no uniqueness assumption on optimal actions.  Arbitrary
stochastic transitions and mixed continuing rows cannot improve this
number.  An optimal implementation uses the chosen initial row and
pure rows for a minimum transversal, updating solely from the reset
report.  Theorem~\ref{thm:r77-reset-gap} gives an attained finite
variational cost lower which is strictly positive above $V_*$ exactly
at the insufficient label budgets.

The cover is an intrinsic obstruction within physical experiments.
For every nonempty finite set system of proper action subsets,
Theorem~\ref{thm:r77-reset-universality} constructs rational reset
data with commuting cost and effect operators with constant success probability and bounded costs whose
optimal memory is its transversal number plus one.  The unrestricted
value and unique uniform initial row are explicit.  Thus the optimal
memory can encode NP-hard covering instances even when the unrestricted
optimal value is simple.  This is a specific reduction; it is not an
assertion that finite-state controller complexity is new in general.
An exact two-dimensional counterexample shows why semidefinite
failure effects alone do not support \eqref{eq:r77-intro-reset}.

\subsection*{Approximation without a Born model of hidden labels}

The earlier physical-dictionary construction gives logarithmic memory
for the one-angle stopped qubit experiment, and a matching leading
constant under a badly-approximable phase assumption.  Its converse
used a density matrix and Born terminal decoder at each hidden label.
An arbitrary classical simulator need not have this representation.

Suppose $J$ causally prepared heralds have target probabilities at
least $p_*>0$.  After herald $n$, a terminal test $m$ is requested.
Its rejection probability is zero at $m=n$ and at least $s>0$ otherwise.
The request is delayed until after the preparation, and the terminal
decoder receives no uncounted copy of the earlier history.
Theorem~\ref{thm:r77-operational-antichain} shows that a simulator
approximating these joint records within total variation $\delta$
must have at least
\[
 a(J)=\min\{m:\tbinom m{\lfloor m/2\rfloor}\ge J\}
\]
labels whenever $\delta\le p_*s^2/(32J^2)$.  The certificate uses only
a finite family of preparation--test experiments.  It starts from joint
herald-and-rejection events, controls herald mass, and only then divides
to obtain a conditional factorization.  A robust support-antichain
argument permits completely arbitrary stochastic decoder columns.

For the original undamped qubit instrument with Bloch rotation
$z=(3+4i)/5$, Theorem~\ref{thm:r77-unrestricted-qubit} gives the
explicit finite certificate
\begin{equation}\label{eq:r77-intro-arithmetic}
 \delta\le\frac{25\zeta}{128\cdot200^J J^2}
 \quad\Longrightarrow\quad N_{\rm cl}(\delta;B)\ge a(J),
 \qquad\zeta=\min\{1,7B/8\}.
\end{equation}
The proof uses rare reset--rotation prefixes, delayed terminal
projectors and an integer separation bound for powers of $z$.
It needs no badly-approximable assertion about the angle.  At fixed
$B>0$ it yields a lower order $\log\log(1/\delta)$ for either start
convention.  An elementary algebraic norm argument extends the same
order to every non-root algebraic phase on the unit circle.
The inherited physical construction is admissible in this larger class
and has order $\log(1/\delta)$; these bounds are not claimed to match.
The sharp antichain example concerns a preparation--test table, not
the qubit's entire chronological dynamics.

\subsection*{Classical ingredients and retained mathematical structure}

Residual equivalence, probabilistic automata and positive realization
are classical \cite{RabinScott59,Rabin63,Heller65,CarlylePaz71,BenvenutiFarina04}.
So are stochastic group structure, representation theory and equivariant
factorization \cite{Schwarz64,Flor69,Serre77,GPT13}.  The shared
initialization invariant and noncyclic joint resource are the additional
statements proved here.  Bellman and stochastic-shortest-path theory
\cite{BertsekasTsitsiklis91}, finite stochastic controller search
\cite{Meuleau99,PoupartBoutilier03}, and finite covering complexity
\cite{Karp72} are used with their established meanings.  Nonnegative
rank and antichain lower bounds have explicit antecedents
\cite{GPT13,VandaeleGillisGlineur17}; the new operational claim includes
the perturbation margin, stopped-cost accounting and algebraic-qubit
application.  A targeted comparison does not certify exhaustive priority.

The general actual-history comparison and its common paid construction
remain in the main article.  Its lower witness is evaluated on the
competitor's actual reached histories.  The main also retains the
general stationary-initialization identity, separately from the new
shared-preparation charge.  The supplement preserves all earlier
finite-order quantum measurements, two-stage reset, memory and
communication, rank-constrained discrimination, performed-law
compactness, variational and dual attainment, optimal strategy
implementation, spectral memory and acquired-geometry arguments.
The independent R58 companion remains separately readable and buildable.
Their scopes are not replaced by a single group example or by the
new positive-effect reset class.  In particular, no result about
response-memory complexity alone solves the original general-$M$
squared-sign minimax-risk frontier or the independent A1/A2 programme.
